\section{Properties of the Lagrangian L}
    The Lagrangian has some very important properties we have to discuss before we look at an example for how to use the Euler-Lagrange-Equations to find the equations of motion for a mechanical system. 
    
    \subsection*{Scale Invariance of the Euler-Lagrange Equations}
    We can multiply the Lagrangian by an arbitrary constant factor without changing the Euler-Lagrange-Equations (Which remember are the equations of motion for the mechanical system).
    Lets proof this. Consider a Lagrange function $L'$ that is scaled by a constant factor $\alpha$
    \begin{align}
    	L'=L\cdot\alpha
    \end{align}
    Lets plug this into our ELE (Euler-Lagrange-Equations) and look what happens.
    
    \begin{align}
        \frac{\partial L'}{\partial q}-\frac{\text{d}}{\text{d}t}\left(\frac{\partial L'}{\partial \dot{q}}\right) = 0\\
        \frac{\partial}{\partial q}\left[L\cdot\alpha\right]-\frac{\text{d}}{\text{d}t}\left(\frac{\partial}{\partial \dot{q}}\left[L\cdot\alpha\right]\right) = 0\\
        \alpha\frac{\partial}{\partial q}\left[L\right]-\alpha\frac{\text{d}}{\text{d}t}\left(\frac{\partial}{\partial \dot{q}}\left[L\right]\right) = 0
    \end{align}
    We can now of course divide by $\alpha$ and obtain
    \begin{align}\frac{\partial}{\partial q}\left[L\right]-\frac{\text{d}}{\text{d}t}\left(\frac{\partial}{\partial \dot{q}}\left[L\right]\right) = 0
    \end{align}
    This now of course is exactly the ELE for L. Therefore scaling $L$ by a constant factor $\alpha$ does not affect the ELE. 
    We will later see, when we work through examples, that we can use this Property of $L$ to simplify our calculations. 
    
    \subsection*{Uniqueness of the Lagrangian up to a Total Derivative}
    We can add the total time-derivative of a function to $L$ without affecting the ELE, i.e.
    \begin{align}
    	L' = L + \frac{\text{d} f}{\text{d}t}
    \end{align}
    
    Now finally, let's look at an example.
